% Pursuing Stacks (cotes 134-2 to 134-8): archive cote 134-4, site page 93 (PDF page 94), typescript page 180.
% typescript. Draft compared: ps3.tex:4767-4832.
% AI-transcribed by Claude Opus 5.5 (claude-opus-5-5, adaptive thinking, effort high), Message Batches API, batch ps2 (msgbatch_01NVwr1Yuo8qJf5rykeGBXPK), 2026-11-12; not yet checked by eye.
\scanpage{134-4/93}
\mspage{180}

9.4.

6\over{5}{6}.     After writing down nicely\hand{,} in the end\hand{,} that long promised key
result,  I thought next thing would be to pull myself up by my bootstraps\unsure{,}
getting the similar result first for test functors $A\longrightarrow B^{\wedge}$ with values
in an elementary modelizer $B^{\wedge}$,  and then for general ``canonical'' modeli-%
zers (M,W).  So I did a little scratchwork pondering along those lines,
before resuming the type writer - engined work.  What then turned out,
it seems,  is that there was'nt any need at all to pass through the parti%
-cular case of M = (Cat) and the somewhat painstaking analysis of our \xout{thre}
three-step diagram on page 96 ,  Finally,  the most useful result of all
the eighty pages grinding\hand{,} since that point\hand{,} is by no means the\xout{t} so-called
key result,  as I anticipated  -  the day after I finally wrote it down,
it was already looking rather ``étriqué'' -  why all this fuss about \xout{(Cat)}
the special case of \xout{(Cat)} test functors with values in (Cat) !  The main
result has been finally more psychological than technical  -  namely
drawing attention,  in the long last,  to the key role of contractible
objects and,  more specifically,  of the contractibility structures asso-%
ciated to modelizers \inserted[typed]{(M,W)}, suggesting that the localizer W \xout{could} should,
conversely,  be describable in terms of the homotopy structure $h_{W}$ .  This
was point d)  in the ``provisional plan of work'' contemplated earlier this
week (page 138)  -  by then I had already the feeling this approach via
d)  would turn out to be the most ``expedient'' one  -  but it was by then
next to impossible for me to keep pushing off still more the approach b),
\xout{which} via test functors with values in (Cat),  which I finally carried
through.  One point which was'nt wholly clear yet that day,  as it is
now,  is the crucial role played by the circumstance that for the really
nice modelizers (surely for those I'm going to call ``canonical'' in the
end),  the associated homotopy structure is indeed a \uline{contractibility}
structure.  Here,  as so often in mathematics (and even outside of mathe-%
matics ...),  the main thing to dig out and discover is where the empha-%
sis belongs  -  which are the \deleted{main} \inserted{really essential} facts or notions or features within a
given context,  and which are accessory\hand{,} namely, \inserted{which} will follow suit by them%
-selves.  It took a while before I would listen to what the things I was
in were insistently telling me.  It finally got through I feel,  and I
believe that from this point \inserted{on} the whole modelizing story is going to go
through extremely smoothly .

     Before starting with the work,  just some retrospective,  somewhat
more technical comments,  afterthoughts rather I should say. \inserted{First of all,} I am not so
happy after all with the \xout{decision} terminological revi\over{\unsure{ew}}{ew} a few days ago

\margin{93}
